\subsection{Proof Construction and Revision}

A proof is a checkable argument that derives the stated conclusion from declared hypotheses, definitions, and earlier proved results.

\begin{proofmethod} \label{meth:proof-construction-audit} \hlt{Proof Construction and Audit}\\
Given an audited theorem-style statement, construct a proof with independently checkable dependencies and obligations.
\begin{enumerate}[label=(\roman*)]
\setlength{\itemsep}{0pt}
\item \hlt{Parse}: identify the objects, hypotheses, quantifier dependencies, and conclusion by Proof Method \ref{meth:proof-grammar}.
\item \hlt{Plan}: select a direct argument, cases, two inclusions, induction, contradiction, construction, or another fitting architecture.
\item \hlt{Dependencies}: record only earlier results whose hypotheses are verified at their citation points.
\item \hlt{Obligations}: assign each numbered claim, implication direction, existence witness, and uniqueness assertion to a proof unit.
\item \hlt{Execution}: bind variables before use and justify every inference by a definition, calculation, hypothesis, or earlier result.
\item \hlt{Boundaries}: handle empty, zero, singular, endpoint, and other degenerate cases that affect an argument.
\item \hlt{Closure}: verify that the final sentence discharges the original conclusion rather than a nearby surrogate.
\item \hlt{Cold read}: test notation, citations, and transitions without relying on memory of the intended argument.
\end{enumerate}
The audit passes only when every obligation is discharged and each dependency is valid at its point of use. When the route is not apparent, expose the plan while retaining every step needed for verification. Mathematical-writing guidance likewise emphasizes status, motivation, and conceptual structure \citep{mit1998mathematicalwriting,knuth1988mathematicalwriting}.
\end{proofmethod}

\begin{process} \label{proc:proof-revision-passes} \hlt{Proof Revision Passes}\\
Revise a completed proof in three ordered passes.
\begin{enumerate}[label=(\roman*)]
\setlength{\itemsep}{0pt}
\item \hlt{Validity}: check every dependency, implication, case split, witness, index, limit, and exceptional case.
\item \hlt{Architecture}: align proof units with statement claims and make long operative stages visible.
\item \hlt{Exposition}: remove repetition, bind notation before use, compress routine algebra, and expand decisive unfamiliar steps.
\end{enumerate}
Review typography only after the mathematical and logical structure is stable.
\end{process}

\begin{warning} \label{warn:evidence-is-not-proof} \hlt{Evidence Is Not Proof}\\
Finite testing, a diagram, numerical agreement, or a successful program run supports only the cases and conditions examined. Such evidence cannot close an unrestricted quantified claim without an exhaustive argument and a proof that the tested cases cover the domain. Computation may discover a conjecture, locate a counterexample, or verify a formally specified finite case; otherwise provide a complete proof or narrow the claim to the verified scope.
\end{warning}

\begin{principle} \label{princ:proof-detail-proportionality} \hlt{Proof Detail Proportionality}\\
Allocate detail where correctness or understanding depends on it. Expand novel constructions, fragile estimates, changes of quantifier, and boundary arguments; compress familiar algebra only while the remaining chain is checkable. The target is a reproducible mathematical argument, neither maximal formality nor persuasive rhetoric.
\end{principle}
